\chapter{定積分と微積分の基本定理}\label{章:定積分}
\section{細かい区間の寄与を足す}\label{節:リーマン和}
\chapterref{章:導入}では，長方形の幅を小さくして面積を求めた。そのときは等しい幅に分け，左端や右端の高さを使った。しかし，同じ図形の面積が，幅の分け方や高さを選ぶ点によって変わっては困る。そこで，まず分け方と点の選び方を明示し，それらに依存しない極限が存在するときに定積分を定義する。

$a<b$ とし，区間 $[a,b]$ を
\[
 a=x_0<x_1<\cdots<x_n=b
\]
と分ける。この点の組を\textbf{分割} $P$ といい，第 $k$ 小区間の幅を $\Delta x_k=x_k-x_{k-1}$ と書く。さらに，各小区間 $[x_{k-1},x_k]$ の中から点 $\xi_k$ を一つずつ選ぶ。$f$ が $[a,b]$ で有界なら，
\[
 S(P,\xi)=\sum_{k=1}^{n}f(\xi_k)\Delta x_k
\]
を作れる。これを\textbf{リーマン和}という。$f(\xi_k)$ が正なら高さと幅の積であり，負なら符号を付けた寄与として足し合わせる。

\begin{figure}[htbp]
\centering
\begin{tikzpicture}[x=3cm,y=1.8cm,>=Stealth]
 \filldraw[fill=blue!12,draw=blue!55] (0,0) rectangle (0.8,1.1225);
 \filldraw[fill=blue!12,draw=blue!55] (0.8,0) rectangle (1.8,1.01);
 \filldraw[fill=blue!12,draw=blue!55] (1.8,0) rectangle (3,1.5625);
 \draw[->] (-0.1,0)--(3.35,0) node[right]{$x$};
 \draw[->] (0,-0.13)--(0,2.4) node[above]{$y$};
 \draw[very thick,domain=0:3.1,samples=80] plot (\x,{1+0.25*(\x-1)^2});
 \foreach \xx/\hh/\name in {0.3/1.1225/1,1.2/1.01/2,2.5/1.5625/3}{
  \draw[blue,dashed] (\xx,0)--(\xx,\hh);
  \fill[blue] (\xx,\hh) circle(1.7pt);
  \node[blue] at (\xx,-0.45) {$\xi_{\name}$};
 }
 \node[below left] at (0,0) {$a=x_0$};
 \node[below] at (0.8,0) {$x_1$};
 \node[below] at (1.8,0) {$x_2$};
 \node[below right] at (3,0) {$x_3=b$};
 \node at (2.7,2.28) {$y=f(x)$};
\end{tikzpicture}
\caption{幅が異なっていてもよく，各小区間のどの点の高さを使ってもよい。}
\label{図:リーマン和}
\end{figure}

図\ref{図:リーマン和}のように，分割を等間隔にする必要はなく，代表点 $\xi_k$ を端点にする必要もない。分割を細かくするとは，単に小区間の個数を増やすことではない。一つだけ幅の広い区間が残っている可能性を除くため，最大の幅
\[
 \|P\|=\max_{1\le k\le n}\Delta x_k
\]
を $0$ に近づける。

\begin{definition}{定積分}
$[a,b]$ 上の有界な関数 $f$ に対し，$\|P\|\to0$ のとき，分割と代表点の選び方によらずリーマン和が同じ有限値 $I$ に近づくなら，$f$ は $[a,b]$ で\textbf{積分可能}であるという。その値を\textbf{定積分}と呼び，
\[
 \int_a^b f(x)\dd x=I
\]
と書く。
\end{definition}

正確には，任意の $\varepsilon>0$ に対して，ある $\delta>0$ を選び，$\|P\|<\delta$ を満たすすべての分割とすべての代表点について $|S(P,\xi)-I|<\varepsilon$ とできることである。特定の分け方で一つの極限が求まるだけでは，この定義を満たすことまでは証明できない。

記号 $\int$ は和を極限まで進める操作を表し，$\dd x$ はどの変数に沿った幅を足しているかを表す。$
\int_a^b f(t)\dd t$ と文字を替えても同じ数である。$x$ や $t$ はこの積分の内部で用いる変数であり，$a,b$ が固定されていれば，定積分全体は一つの数になる。

\section{連続な関数は積分できる}\label{節:積分可能性}
定義は積分が存在する場合の意味を定めたが，存在そのものをまだ保証していない。本書で主に扱う連続関数については，次の定理を用いる。

\begin{theorem}{連続関数の積分可能性：採用する事実}\label{定理:連続積分可能}
閉区間 $[a,b]$ で連続な関数は，その区間でリーマン積分可能である。
\end{theorem}

この定理の完全な証明は，本書では扱わない。考え方は，十分短い小区間では関数の高さの変動も小さくなり，どの代表点を選んでも和の差が小さくなるというものである。閉区間上の連続性を使うと，その小ささを区間全体でそろえられる。各小区間での最大の高さと最小の高さの差が $\eta$ 以下なら，上から作る和と下から作る和の差は
\[
 \sum_{k=1}^{n}\eta\Delta x_k=\eta(b-a)
\]
以下になる。このような評価をすべての分割について保証する部分も含め，上の定理として採用する。

したがって，連続関数の定積分の値を計算するときは，存在をこの定理で確かめたうえで，計算しやすい等分割と代表点を選んでよい。右端を選ぶなら，
\[
 \int_a^b f(x)\dd x
 =\lim_{n\to\infty}\frac{b-a}{n}
  \sum_{k=1}^{n}f\left(a+\frac{k(b-a)}n\right)
\]
となる。ここでは幅が $\frac{b-a}{n}$，第 $k$ 小区間の右端が $a+\frac{k(b-a)}n$ である。

\begin{example}{リーマン和から値を求める}
区間 $[1,3]$ を $n$ 等分し，各小区間の右端を使って $\int_1^3(2x-1)\dd x$ を求めよ。
\end{example}
被積分関数は連続なので，定積分は存在する。幅は $\frac2n$，右端は $1+\frac{2k}{n}$ である。したがって，リーマン和は
\begin{align*}
 S_n
 &=\frac2n\sum_{k=1}^{n}\left\{2\left(1+\frac{2k}n\right)-1\right\}\\
 &=\frac2n\sum_{k=1}^{n}\left(1+\frac{4k}n\right)\\
 &=2+\frac8{n^2}\cdot\frac{n(n+1)}2
  =6+\frac4n.
\end{align*}
よって $n\to\infty$ として，定積分は $6$ となる。この例では和の公式が使えたが，一般の関数では和を直接計算するのは難しい。後で微分との関係を使い，別の計算法を導く。

\section{定積分の基本性質}\label{節:定積分性質}
以下では，特に断らない限り，被積分関数は扱う閉区間で連続であるとする。まず，積分の端点が同じ場合と，順序が逆の場合を
\[
 \int_a^a f(x)\dd x=0,\qquad
 \int_b^a f(x)\dd x=-\int_a^b f(x)\dd x\quad(a<b)
\]
と定める。積分の向きが反対になると符号が反対になる。この約束により，区間をつなぐ式を端点の大小にかかわらず統一して書ける。

\begin{theorem}{線形性と区間の加法性}\label{定理:積分基本性質}
$\alpha,\beta$ を実数とする。関係する区間で $f,g$ が連続なら，
\begin{align*}
 \int_a^b\{\alpha f(x)+\beta g(x)\}\dd x
 &=\alpha\int_a^b f(x)\dd x+\beta\int_a^b g(x)\dd x,\\
 \int_a^b f(x)\dd x
 &=\int_a^c f(x)\dd x+\int_c^b f(x)\dd x
\end{align*}
が成り立つ。後者は $a,c,b$ の順序によらない。
\end{theorem}

線形性は，同じ分割と代表点を使った有限和について
\[
 \sum_k\{\alpha f(\xi_k)+\beta g(\xi_k)\}\Delta x_k
 =\alpha\sum_k f(\xi_k)\Delta x_k
  +\beta\sum_k g(\xi_k)\Delta x_k
\]
と分け，極限をとれば得られる。区間の加法性は，$a<c<b$ のとき，分割に $c$ を含めて和を二つに分ければ分かる。端点の順序が異なる場合は，逆向きの積分の定義から同じ式になる。積分区間を分けても，被積分関数を勝手に変えているわけではない。

\begin{theorem}{大小関係と絶対値の評価}\label{定理:積分評価}
$a\le b$ とする。区間 $[a,b]$ で連続な関数について，次が成り立つ。
\begin{align*}
 f(x)\le g(x)\ (a\le x\le b)
 &\quad\Longrightarrow\quad
 \int_a^b f(x)\dd x\le\int_a^b g(x)\dd x,\\
 \left|\int_a^b f(x)\dd x\right|
 &\le\int_a^b|f(x)|\dd x.
\end{align*}
特に，$m\le f(x)\le M$ なら，
\[
 m(b-a)\le\int_a^b f(x)\dd x\le M(b-a).
\]
\end{theorem}

最初の性質は，各小区間の幅が正であり，リーマン和の各項について大小関係が保たれることから分かる。二つ目は $-|f|\le f\le|f|$ を積分すれば得られる。最後の評価では，定数関数の積分が高さと区間の長さの積になることを使った。これらの不等式は，端点を $a\le b$ の順に置いていることが前提である。

\begin{example}{原始関数を探さずに範囲を決める}
$\int_0^1\frac1{1+x^2}\dd x$ が $\frac12$ 以上 $1$ 以下であることを示せ。
\end{example}
$0\le x\le1$ では $1\le1+x^2\le2$ であり，すべて正なので逆数の大小は逆になる。
\[
 \frac12\le\frac1{1+x^2}\le1.
\]
区間の長さは $1$ だから，これを積分して
\[
 \frac12\le\int_0^1\frac1{1+x^2}\dd x\le1
\]
を得る。積分の値を正確に計算できなくても，関数の高さを比較すれば範囲を保証できる。この例の原始関数は\chapterref{章:置換の背景}で扱う。

\subsection{関数の平均値}
$a<b$ のとき，
\[
 f_{\mathrm{av}}=\frac1{b-a}\int_a^b f(x)\dd x
\]
を $f$ の区間 $[a,b]$ における\textbf{平均値}という。等分割のリーマン和を区間の長さで割ると $\frac1n\sum_{k=1}^{n}f(\xi_k)$ になるので，有限個の値の平均を細かくしていったものと解釈できる。微分で用いた平均変化率とは別の量である。

$f$ の最小値を $m$，最大値を $M$ とすると，積分の評価から $m\le f_{\mathrm{av}}\le M$ である。連続関数は最小値と最大値の間の値をすべてとるから，ある $c\in[a,b]$ で
\[
 \int_a^b f(x)\dd x=f(c)(b-a)
\]
が成り立つ。これを\textbf{積分の平均値の定理}という。高さを一定の $f(c)$ にした長方形と，元の関数の定積分が等しくなるという意味である。

\section{符号付きの積分と面積を区別する}\label{節:積分符号面積}
$f\ge0$ のとき，リーマン和は長方形の面積の和である。この極限として，曲線の下の面積を定める。連続関数が負になる部分では，定積分への寄与は負になるため，定積分をそのまま幾何学的な面積と呼ぶことはできない。

$f$ が正の部分の面積を $S_+$，負の部分と $x$ 軸の間の面積を $S_-$ とすれば，
\[
 \int_a^b f(x)\dd x=S_+-S_-,\qquad
 \int_a^b|f(x)|\dd x=S_++S_-
\]
である。どちらも同じグラフから求めるが，符号を保つか，面積としてすべて正にするかが異なる。面積を求めるには零点を調べ，符号が変わる場所で区間を分ける。

\begin{example}{同じグラフから二つの量を求める}
$y=x-1$ の $0\le x\le3$ における定積分と，グラフと $x$ 軸の間の面積を，三角形の面積から求めよ。
\end{example}
グラフは $x=1$ で $x$ 軸と交わる。$0\le x\le1$ の下側の三角形の面積は $\frac12$，$1\le x\le3$ の上側の三角形の面積は $2$ である。したがって，
\[
 \int_0^3(x-1)\dd x=-\frac12+2=\frac32,
 \qquad
 \int_0^3|x-1|\dd x=\frac12+2=\frac52.
\]
図\ref{図:定積分符号}の下側の部分は，左の計算では引き，右の計算では足している。

\begin{figure}[htbp]
\centering
\begin{tikzpicture}[x=2.3cm,y=1.5cm,>=Stealth]
 \fill[red!15] (0,0)--(0,-1)--(1,0)--cycle;
 \fill[blue!15] (1,0)--(3,2)--(3,0)--cycle;
 \draw[->] (-0.2,0)--(3.45,0) node[right]{$x$};
 \draw[->] (0,-1.3)--(0,2.5) node[above]{$y$};
 \draw[very thick,domain=0:3.2] plot (\x,{\x-1});
 \draw[dashed] (3,0)--(3,2)--(0,2);
 \node[below] at (1,0) {$1$};
 \node[below] at (3,0) {$3$};
 \node[left] at (0,-1) {$-1$};
 \node[left] at (0,2) {$2$};
 \node[red!70!black] at (0.3,-0.32) {$S_-$};
 \node[blue!70!black] at (2.4,0.48) {$S_+$};
 \node at (2.45,2.25) {$y=x-1$};
\end{tikzpicture}
\caption{定積分は上側の面積から下側の面積を引く。幾何学的な面積は両方を足す。}
\label{図:定積分符号}
\end{figure}

\section{積分の上端を動かして微分する}\label{節:基本定理第一}
ここまで，定積分をリーマン和の極限として定義し，和から分かる性質を調べた。次に，右端を変数にして積分そのものを関数として見る。$f$ が $[a,b]$ で連続なとき，
\[
 A(x)=\int_a^x f(t)\dd t\qquad(a\le x\le b)
\]
とおく。$x$ が動くと，上端までに足し合わせる範囲が変わる。$f\ge0$ なら面積を表すが，負の値を含む場合にも符号付きの量として同じ式を使う。

\begin{figure}[htbp]
\centering
\begin{tikzpicture}[x=2.5cm,y=1.5cm,>=Stealth]
 \fill[blue!12] (0,0)--plot[domain=0:2,samples=60] (\x,{1+0.2*\x*\x})--(2,0)--cycle;
 \fill[red!18] (2,0)--plot[domain=2:2.5,samples=35] (\x,{1+0.2*\x*\x})--(2.5,0)--cycle;
 \draw[->] (-0.15,0)--(3.15,0) node[right]{$t$};
 \draw[->] (0,-0.1)--(0,3.15) node[above]{$y$};
 \draw[very thick,domain=0:3,samples=80] plot (\x,{1+0.2*\x*\x});
 \draw[dashed] (2,0)--(2,1.8);
 \draw[dashed] (2.5,0)--(2.5,2.25);
 \node[below] at (0,0) {$a$};
 \node[below] at (2,0) {$x$};
 \node[below] at (2.5,0) {$x+h$};
 \node[blue!70!black] at (1,0.55) {$A(x)$};
 \node[red!70!black] at (2.25,0.8) {増分};
 \node at (2.7,3.05) {$y=f(t)$};
\end{tikzpicture}
\caption{上端を $h$ だけ動かしたときの積分の増分は，幅 $h$ の区間の積分になる。}
\label{図:面積関数増分}
\end{figure}

図\ref{図:面積関数増分}は $h>0$ で $f$ が正の場合を描いている。幅が小さいとき，増分は高さ $f(x)$，幅 $h$ の長方形に近い。これを符号や $h$ の向きに依存しない形で示す。

\begin{theorem}{微積分の基本定理：積分を微分する}\label{定理:基本定理第一}
$f$ が $[a,b]$ で連続なら，$A(x)=\int_a^x f(t)\dd t$ は $[a,b]$ で連続であり，内点で微分可能で，
\[
 A'(x)=f(x)\qquad(a<x<b)
\]
を満たす。端点では，区間の内側からの片側微分で同じ式が成り立つ。
\end{theorem}

まず，連続関数 $f$ は閉区間で有界だから，$|f(t)|\le M$ となる $M$ がある。積分の評価と区間の加法性により，任意の $u,v\in[a,b]$ について
\[
 |A(v)-A(u)|=\left|\int_u^v f(t)\dd t\right|
 \le M|v-u|
\]
となる。よって $A$ は連続である。

微分を調べるため，$x,x+h\in[a,b]$，$h\ne0$ とする。区間の加法性から
\begin{align*}
 \frac{A(x+h)-A(x)}h-f(x)
 &=\frac1h\int_x^{x+h}f(t)\dd t-f(x)\\
 &=\frac1h\int_x^{x+h}\{f(t)-f(x)\}\dd t.
\end{align*}
ここで，定数関数 $f(x)$ の $x$ から $x+h$ までの積分が $hf(x)$ になることを使った。

任意の $\varepsilon>0$ に対し，$f$ の $x$ での連続性から，$|t-x|<\delta$ なら $|f(t)-f(x)|<\varepsilon$ となる $\delta>0$ を選べる。$0<|h|<\delta$ のとき，$x$ と $x+h$ の間のすべての $t$ がこの条件を満たす。積分区間の長さは $|h|$ だから，向きがどちらであっても
\[
 \left|\frac1h\int_x^{x+h}\{f(t)-f(x)\}\dd t\right|
 \le\frac1{|h|}\,\varepsilon|h|=\varepsilon.
\]
したがって差商は $f(x)$ に近づき，$A'(x)=f(x)$ を得る。図から予想した「積分の増加率が高さになる」という関係が，負の高さや $h<0$ の場合も含めて証明できた。

\section{原始関数から定積分を計算する}\label{節:原始関数基本定理}
基本定理は，積分によって作った関数を微分すると元に戻ることを示した。今度は，微分すると $f$ になる関数を先に一つ見つけ，それを積分の計算に使う。

\begin{definition}{原始関数と不定積分}
一つの区間で $F'(x)=f(x)$ を満たす関数 $F$ を，$f$ の\textbf{原始関数}という。すべての原始関数を表す記法として
\[
 \int f(x)\dd x=F(x)+C
\]
と書き，\textbf{不定積分}と呼ぶ。$C$ は任意の定数であり，\textbf{積分定数}という。
\end{definition}

なぜ定数を足すだけで，すべての原始関数を表せるのだろうか。$F,G$ がともに $f$ の原始関数なら，$(F-G)'=0$ である。前章の定理\ref{定理:導関数同じ}により，一つの区間では $F-G$ が定数になるからである。逆に，$F+C$ を微分すれば $F'=f$ となる。たとえば $x^2$ の原始関数は $\frac{x^3}{3}+C$ である。

ここでも区間ごとに考える必要がある。$\frac1x$ の定義域は $x<0$ と $x>0$ に分かれるので，それぞれの区間で $\log|x|+C$ と書けるが，二つの区間で同じ定数を使う必要はない。また，基本定理から連続関数には原始関数が存在するが，その原始関数が既知の関数の式で簡単に表せるとは限らない。

\begin{theorem}{微積分の基本定理：原始関数で計算する}\label{定理:基本定理第二}
$f$ が $[a,b]$ で連続で，$F$ が $[a,b]$ で連続，$(a,b)$ で $F'=f$ を満たすなら，
\[
 \int_a^b f(x)\dd x=F(b)-F(a)
\]
が成り立つ。右辺を $[F(x)]_a^b$ とも書く。
\end{theorem}

証明には $A(x)=\int_a^x f(t)\dd t$ を使う。基本定理の前半から $A'=f$ なので，$(F-A)'=0$ である。$F-A$ は閉区間で連続だから，平均値の定理により端点を含む区間全体で一定となる。$A(a)=0$ を用いれば
\[
 F(x)-A(x)=F(a),\qquad A(x)=F(x)-F(a).
\]
ここに $x=b$ を代入すると，求める公式を得る。逆向きの積分や両端が同じ場合も，それぞれの定義から同じ計算式になる。

この証明の順序を確かめておこう。定積分は，初めから原始関数の差として定義したのではない。リーマン和の極限として定義し，その積分を微分すると $f$ に戻ることを示し，最後に平均値の定理によって任意の原始関数と結び付けた。これにより，長方形の和で表される量を，微分の逆を探す計算で求められるようになった。

\begin{example}{微分して確かめてから端点を代入する}
次の定積分を求めよ。
\[
 \text{(1)}\quad\int_0^2(3x^2-4x+1)\dd x,\qquad
 \text{(2)}\quad\int_1^{\e}\frac1x\dd x,\qquad
 \text{(3)}\quad\int_0^{\frac\pi2}\cos x\dd x.
\]
\end{example}
(1) では，微分すると $3x^2-4x+1$ になる関数を項ごとに探す。$F(x)=x^3-2x^2+x$ とすれば $F'=3x^2-4x+1$ だから，
\[
 \int_0^2(3x^2-4x+1)\dd x
 =[x^3-2x^2+x]_0^2=(8-8+2)-0=2.
\]
不定積分では $F+C$ と書くが，定積分では $(F(b)+C)-(F(a)+C)$ の定数が打ち消し合う。計算結果に積分定数を付ける必要はない。

(2) は区間全体で $x>0$ だから，$\log x$ を原始関数として使える。
\[
 \int_1^{\e}\frac1x\dd x=[\log x]_1^{\e}=1-0=1.
\]
(3) では $(\sin x)'=\cos x$ より，
\[
 \int_0^{\frac\pi2}\cos x\dd x
 =[\sin x]_0^{\frac\pi2}=1-0=1.
\]
また，最初の例題で和から求めた積分は，$\int_1^3(2x-1)\dd x=[x^2-x]_1^3=6$ と短く計算できる。二つの方法が同じ値になる理由が，基本定理で保証されている。

\section{上下端が変わる積分}\label{節:変数上下端}
$\int_a^x f(t)\dd t$ の微分が分かれば，上端や下端が別の関数になっている場合も合成関数の微分で扱える。$f$ がある区間で連続で，$u(x),v(x)$ が微分可能，その値がこの区間内にあるとする。区間内の固定した点 $c$ を選び，$A(s)=\int_c^s f(t)\dd t$ とおけば，
\[
 \int_{u(x)}^{v(x)}f(t)\dd t=A(v(x))-A(u(x)).
\]
両辺を微分し，$A'=f$ を使うと，
\[
 \boxed{\frac{\dd}{\dd x}\int_{u(x)}^{v(x)}f(t)\dd t
 =f(v(x))v'(x)-f(u(x))u'(x)}
\]
を得る。下端からの寄与に負号が付くのは，区間の加法性でその部分を引いているからである。$u(x)$ と $v(x)$ の大小が途中で入れ替わっても，向き付きの積分の約束によりこの式を使える。

\begin{example}{原始関数の式を求めずに微分する}
$G(x)=\int_x^{x^2}\e^{t^2}\dd t$ を微分せよ。
\end{example}
被積分関数 $\e^{t^2}$ はすべての実数で連続である。上端 $x^2$ の微分は $2x$，下端 $x$ の微分は $1$ だから，
\[
 G'(x)=\e^{(x^2)^2}\cdot2x-\e^{x^2}\cdot1
       =2x\e^{x^4}-\e^{x^2}.
\]
上端を代入してから内側の関数の微分を掛けるので，最初の指数は $x^4$ になる。$\e^{t^2}$ の原始関数を具体的な式で表す必要はない。

ただし，この公式では被積分関数を $f(t)$ とし，積分の外側の変数 $x$ がその中には現れないとしている。たとえば
\[
 H(x)=\int_0^x xt\dd t
      =x\int_0^x t\dd t=\frac{x^3}{2}
\]
では，$H'(x)=\frac{3x^2}{2}$ となる。上端を代入した $x^2$ だけでは足りない。積分の間は $x$ を定数として外に出せても，積分全体を $x$ で微分するときには，その因子も変わるからである。

\section{積分を含む等式から関数を決める}\label{節:積分方程式}
未知の関数が積分の中にも現れる等式を，\textbf{積分方程式}という。上端が変数なら，微分によって積分を外せる場合がある。ただし，微分は定数の情報を失うため，元の等式から得られる値の条件を残し，最後に元の式へ代入して確かめる。

\begin{example}{微分と初期値から関数を求める}
連続関数 $f$ が，すべての実数 $x$ について
\[
 f(x)=2+\int_0^x f(t)\dd t
\]
を満たすとき，$f$ を求めよ。
\end{example}
$x=0$ を代入すると，まず $f(0)=2$ と分かる。$f$ が連続なので，右辺の積分は基本定理により微分可能である。したがって左辺の $f$ も微分可能となり，両辺を微分して $f'(x)=f(x)$ を得る。

この等式だけで直ちに $f=\e^x$ としてはいけない。定数倍も同じ微分の関係を満たすからである。$\e^{-x}$ を掛けた関数を微分すると，
\[
 \frac{\dd}{\dd x}\{\e^{-x}f(x)\}
 =\e^{-x}\{f'(x)-f(x)\}=0
\]
となる。導関数が $0$ なので $\e^{-x}f(x)$ は定数であり，$x=0$ の値からその定数は $2$ である。よって候補は $f(x)=2\e^x$ となる。

最後に元の式を確認する。
\[
 2+\int_0^x2\e^t\dd t
 =2+[2\e^t]_0^x=2+2\e^x-2=2\e^x.
\]
確かに成り立つので，求める関数は $f(x)=2\e^x$ である。導いた条件が必要条件であるだけでなく，元の積分方程式を満たすことまで確かめた。

ここまでで，定積分の意味と，その計算を微分に結び付ける根拠がそろった。次章からは，原始関数を見つけるための式の整理を詳しく扱う。どの公式を使う場合も，求めた原始関数を微分して被積分関数に戻るかを確認できる。複雑な計算ほど，この戻り道が役に立つ。

\section*{章末問題}
\begin{enumerate}[label=\textbf{\thechapter.\arabic*},leftmargin=*]
 \item\label{問:定積分リーマン和} 区間 $[0,2]$ を $n$ 等分し，各小区間の右端を使ったリーマン和から $\int_0^2x^2\dd x$ を求めよ。幅と代表点を明記せよ。
 \item\label{問:定積分性質} 連続関数 $f,g$ が
 \[
  \int_0^1f(x)\dd x=2,\qquad
  \int_1^3f(x)\dd x=-1,\qquad
  \int_0^3g(x)\dd x=4
 \]
 を満たす。$\int_0^3f(x)\dd x$ と $\int_3^0\{2f(x)-3g(x)\}\dd x$ を求めよ。
 \item\label{問:定積分と面積} $\int_{-1}^2(x-1)\dd x$ と，直線 $y=x-1$，$x$ 軸，直線 $x=-1,x=2$ で囲まれる部分の面積を求めよ。符号が変わる点で区間を分けて説明せよ。
 \item\label{問:基本定理計算} 基本定理を使って次を求めよ。
 \[
  \text{(1)}\quad\int_{-1}^2(x^2+2x)\dd x,\qquad
  \text{(2)}\quad\int_0^{\frac\pi2}(\sin x+\cos x)\dd x,
 \]
 \[
  \text{(3)}\quad\int_1^{\e}\left(\frac2x+x\right)\dd x.
 \]
 \item\label{問:面積関数増減} $A(x)=\int_1^x(t-2)(t+1)\dd t$ とする。$A'(x)$ を求め，区間 $[-2,3]$ における $A$ の最大値・最小値と，それらをとる $x$ を求めよ。
 \item\label{問:変数上下端} $G(x)=\int_{\cos x}^{\sin x}\sqrt{1+t^2}\dd t$ を微分せよ。下端から生じる項の符号に注意せよ。
 \item\label{問:積分の評価} $0\le x\le1$ で $1\le\sqrt{1+x^2}\le1+\frac{x^2}{2}$ が成り立つことを示し，それを用いて
 \[
  1\le\int_0^1\sqrt{1+x^2}\dd x\le\frac76
 \]
 を証明せよ。平方根の原始関数を求める必要はない。
 \item\label{問:積分方程式基本} 次の条件を満たす連続関数を求めよ。存在しない場合は，その理由を示せ。
 \begin{enumerate}[label=(\arabic*),leftmargin=*]
  \item すべての実数 $x$ で $f(x)=3-\int_0^x f(t)\dd t$。
  \item すべての $x\in[0,1]$ で $g(x)=x+\int_0^1g(t)\dd t$。
 \end{enumerate}
\end{enumerate}

\smallskip\noindent 解答は\pageref{解答:第3章}ページ。
